\documentclass[a4paper,12pt]{article}

\usepackage[T2A]{fontenc}
\usepackage[utf8]{inputenc}
\usepackage[russian]{babel}

\usepackage{amsmath, amssymb, amsthm, mathtools}
\usepackage{helvet}
\renewcommand{\familydefault}{\sfdefault}
\usepackage{icomma}
\usepackage{geometry}
\usepackage{microtype}
\usepackage{tikz}
\usepackage{tkz-euclide}
\usepackage{pgfplots}
\pgfplotsset{compat=1.18}
\usepackage{tabularx, booktabs, longtable}
\usepackage{enumitem}
\usepackage{hyperref}
\usepackage{parskip}
\usepackage{fancyhdr}
\usepackage{setspace}
\usepackage{xcolor}

\geometry{left=18mm,right=18mm,top=18mm,bottom=20mm}

\definecolor{accent}{HTML}{008080}
\definecolor{darkaccent}{HTML}{004D4D}
\definecolor{textgray}{HTML}{333333}
\definecolor{linegray}{HTML}{B8C8C8}
\definecolor{softgray}{HTML}{F3F7F7}

\hypersetup{hidelinks}
\onehalfspacing
\setlength{\parindent}{0pt}
\setlength{\parskip}{6pt}

\pagestyle{fancy}
\fancyhf{}
\lhead{\textcolor{darkaccent}{Тригонометрические уравнения}}
\rhead{\textcolor{darkaccent}{Тимофей Васильченко}}
\cfoot{\thepage}
\renewcommand{\headrulewidth}{0.4pt}
\renewcommand{\headrule}{\hbox to\headwidth{\color{linegray}\leaders\hrule height \headrulewidth\hfill}}

\setlist[itemize]{leftmargin=1.15cm,itemsep=2pt,topsep=3pt}
\setlist[enumerate]{leftmargin=1.1cm,itemsep=3pt,topsep=3pt}

\newcommand{\sectionline}{\vspace{2mm}{\color{linegray}\hrule height 0.5pt}\vspace{2mm}}
\newcommand{\important}[1]{\textcolor{darkaccent}{\textbf{#1}}}
\newcommand{\R}{\mathbb{R}}
\newcommand{\Z}{\mathbb{Z}}

\begin{document}

\begin{titlepage}
\thispagestyle{empty}
\begin{center}
\vspace*{28mm}

{\Large \textcolor{darkaccent}{\textbf{Чек-лист}}}

\vspace{6mm}

{\huge \textbf{Тригонометрические уравнения}}\\[2mm]
{\Large \textbf{формулы преобразования и специальные замены}}

\vspace{18mm}

{\large Лёвин Артём Александрович}\\[1mm]
{\large эксклюзивно для Тимофея Васильченко}

\vspace{12mm}

{\large \today}

\vfill

\begin{tikzpicture}[x=1cm,y=1cm]
\draw[accent,line width=1.1pt]
  (-5.7,0) .. controls (-3.8,0.9) and (-2.1,-0.8) .. (-0.3,0.15)
  .. controls (1.6,1.05) and (3.5,-0.65) .. (5.7,0.2);
\draw[linegray,line width=0.6pt]
  (-5.2,-0.35) .. controls (-3.0,0.25) and (-1.4,-0.5) .. (0.4,-0.2)
  .. controls (2.4,0.2) and (3.7,-0.3) .. (5.2,-0.1);
\end{tikzpicture}

\vspace{12mm}
\end{center}
\end{titlepage}

\setcounter{page}{2}

\section*{\textcolor{darkaccent}{1. Главная цель преобразований}}

В тригонометрических уравнениях почти всегда нужно добиться двух вещей:

\[
\boxed{\text{одинаковые углы}}
\qquad\text{и}\qquad
\boxed{\text{одинаковые функции}}.
\]

Если в уравнении стоят разные углы, например
\[
\cos x,\qquad \cos 2x,\qquad \cos 5x,
\]
то обычное раскрытие двойного угла может не помочь. Тогда нужно искать формулы, которые меняют сумму, разность или произведение тригонометрических функций.

\important{Главная мысль.} Если стандартная формула не приводит к одинаковым углам, ищем специальную формулу.

\sectionline

\section*{\textcolor{darkaccent}{2. Алгоритм поиска идеи в новой строке}}

Каждую новую строку решения удобно проверять по одному и тому же порядку.

\begin{enumerate}
\item Можно ли что-то \important{вынести}?
\item Можно ли применить \important{формулы сокращённого умножения}?
\item Если есть дроби: можно ли привести к \important{общему знаменателю}?
\item Можно ли сделать \important{замену}?
\item Если слагаемых четыре: возможны ли \important{группировка и двойное вынесение}?
\item Если ничего не сработало: ищем \important{специальную тригонометрическую формулу}.
\item Если и это не сработало: думаем о \important{специальном методе}: ограниченность, монотонность, область значений, нестандартная замена.
\end{enumerate}

\important{Важно.} Анализируем не только всё задание целиком, но и каждую новую строку решения.

\sectionline

\section*{\textcolor{darkaccent}{3. Формулы суммы и разности}}

Эти формулы нужны, когда в уравнении встречаются суммы или разности синусов и косинусов с разными углами.

\[
\sin \alpha+\sin \beta
=
2\sin\frac{\alpha+\beta}{2}\cos\frac{\alpha-\beta}{2}.
\]

\[
\sin \alpha-\sin \beta
=
2\cos\frac{\alpha+\beta}{2}\sin\frac{\alpha-\beta}{2}.
\]

\[
\cos \alpha+\cos \beta
=
2\cos\frac{\alpha+\beta}{2}\cos\frac{\alpha-\beta}{2}.
\]

\[
\cos \alpha-\cos \beta
=
-2\sin\frac{\alpha+\beta}{2}\sin\frac{\alpha-\beta}{2}.
\]

\important{Как выбирать пару углов.} Надо пробовать такие пары, после которых появится угол, уже имеющийся в уравнении.

\subsection*{\textcolor{darkaccent}{Пример}}

Решим уравнение:
\[
\cos x+\cos 5x+\cos 2x=0.
\]

Сложим первый и второй косинусы:
\[
\cos x+\cos 5x
=
2\cos 3x\cos(-2x).
\]

Так как $\cos(-2x)=\cos 2x$, получаем:
\[
2\cos 3x\cos 2x+\cos 2x=0.
\]

Выносим общий множитель:
\[
\cos 2x(2\cos 3x+1)=0.
\]

Отсюда:
\[
\cos 2x=0
\qquad\text{или}\qquad
2\cos 3x+1=0.
\]

\[
2x=\frac{\pi}{2}+\pi k
\quad\Longrightarrow\quad
x=\frac{\pi}{4}+\frac{\pi k}{2},
\quad k\in\Z.
\]

\[
\cos 3x=-\frac12
\quad\Longrightarrow\quad
3x=\pm\frac{2\pi}{3}+2\pi k.
\]

\[
x=\pm\frac{2\pi}{9}+\frac{2\pi k}{3},
\quad k\in\Z.
\]

Ответ:
\[
\boxed{
x=\frac{\pi}{4}+\frac{\pi k}{2}
\quad\text{или}\quad
x=\pm\frac{2\pi}{9}+\frac{2\pi k}{3},
\quad k\in\Z
}.
\]

\sectionline

\section*{\textcolor{darkaccent}{4. Формулы произведения}}

Эти формулы применяются, когда в уравнении есть произведение тригонометрических функций.

\[
\sin \alpha\cos \beta
=
\frac12\left(\sin(\alpha+\beta)+\sin(\alpha-\beta)\right).
\]

\[
\cos \alpha\cos \beta
=
\frac12\left(\cos(\alpha+\beta)+\cos(\alpha-\beta)\right).
\]

\[
\sin \alpha\sin \beta
=
\frac12\left(\cos(\alpha-\beta)-\cos(\alpha+\beta)\right).
\]

\important{Зачем это нужно.} Произведение может превратиться в сумму, где углы станут удобнее для сравнения, вынесения или замены.

\subsection*{\textcolor{darkaccent}{Пример}}

\[
2\sin 3x\cos x-\sin 4x=0.
\]

Применяем формулу произведения:
\[
2\sin 3x\cos x
=
\sin 4x+\sin 2x.
\]

Тогда:
\[
\sin 4x+\sin 2x-\sin 4x=0.
\]

\[
\sin 2x=0.
\]

\[
2x=\pi k
\quad\Longrightarrow\quad
\boxed{x=\frac{\pi k}{2},\quad k\in\Z}.
\]

\sectionline

\section*{\textcolor{darkaccent}{5. Формулы понижения степени}}

Формулы понижения степени нужны, когда встречаются квадраты или четвёртые степени.

\[
\sin^2 x=\frac{1-\cos 2x}{2},
\qquad
\cos^2 x=\frac{1+\cos 2x}{2}.
\]

\[
\sin^4 x=\left(\sin^2 x\right)^2
=
\left(\frac{1-\cos 2x}{2}\right)^2.
\]

\[
\cos^4 x=\left(\cos^2 x\right)^2
=
\left(\frac{1+\cos 2x}{2}\right)^2.
\]

\important{При четвёртой степени} сначала представляем её как квадрат квадрата:
\[
\sin^4 x=(\sin^2 x)^2.
\]

Затем применяем формулу понижения степени.

\subsection*{\textcolor{darkaccent}{Пример}}

\[
\sin^4 x+\sin^4\left(x+\frac{\pi}{4}\right)+\sin^4\left(x-\frac{\pi}{4}\right)=1.
\]

Преобразуем каждое слагаемое.

\[
\sin^4 x
=
\left(\frac{1-\cos 2x}{2}\right)^2
=
\frac{1-2\cos 2x+\cos^2 2x}{4}.
\]

\[
\sin^4\left(x+\frac{\pi}{4}\right)
=
\left(
\frac{1-\cos\left(2x+\frac{\pi}{2}\right)}{2}
\right)^2.
\]

Так как
\[
\cos\left(2x+\frac{\pi}{2}\right)=-\sin 2x,
\]
то
\[
\sin^4\left(x+\frac{\pi}{4}\right)
=
\frac{1+2\sin 2x+\sin^2 2x}{4}.
\]

\[
\sin^4\left(x-\frac{\pi}{4}\right)
=
\left(
\frac{1-\cos\left(2x-\frac{\pi}{2}\right)}{2}
\right)^2.
\]

Так как
\[
\cos\left(2x-\frac{\pi}{2}\right)=\sin 2x,
\]
то
\[
\sin^4\left(x-\frac{\pi}{4}\right)
=
\frac{1-2\sin 2x+\sin^2 2x}{4}.
\]

После сложения:
\[
\frac{3-2\cos 2x+\cos^2 2x+2\sin^2 2x}{4}=1.
\]

Домножаем на $4$:
\[
3-2\cos 2x+\cos^2 2x+2\sin^2 2x=4.
\]

Используем
\[
\sin^2 2x=1-\cos^2 2x.
\]

Получаем уравнение только через $\cos 2x$:
\[
3-2\cos 2x+\cos^2 2x+2(1-\cos^2 2x)=4.
\]

\[
\cos^2 2x+2\cos 2x-1=0.
\]

Дальше выполняется замена:
\[
t=\cos 2x,\qquad -1\le t\le1.
\]

\important{Главное.} После понижения степени задача часто сводится к квадратному уравнению относительно одной тригонометрической функции.

\sectionline

\section*{\textcolor{darkaccent}{6. Почему нельзя делить на переменную}}

В уравнении нельзя просто делить на выражение, которое может быть равно нулю.

Например, если получено:
\[
\cos 2x(2\cos 3x+1)=0,
\]
нельзя делить на $\cos 2x$.

Правильно:
\[
\cos 2x=0
\qquad\text{или}\qquad
2\cos 3x+1=0.
\]

\important{Причина.} Если $\cos 2x=0$, то деление на $\cos 2x$ уничтожит целую часть решений.

\sectionline

\section*{\textcolor{darkaccent}{7. Нестандартная замена \texorpdfstring{$\sin x+\cos x$}{sin x + cos x}}}

Если в уравнении одновременно встречаются
\[
\sin x+\cos x
\qquad\text{и}\qquad
\sin x\cos x,
\]
может помочь замена:
\[
t=\sin x+\cos x.
\]

Тогда:
\[
t^2=(\sin x+\cos x)^2.
\]

\[
t^2=\sin^2 x+2\sin x\cos x+\cos^2 x.
\]

Так как
\[
\sin^2 x+\cos^2 x=1,
\]
получаем:
\[
t^2=1+2\sin x\cos x.
\]

Отсюда:
\[
2\sin x\cos x=t^2-1.
\]

\important{Область значений замены:}
\[
-\sqrt2\le \sin x+\cos x\le \sqrt2.
\]

Значит:
\[
t\in[-\sqrt2;\sqrt2].
\]

\subsection*{\textcolor{darkaccent}{Пример}}

\[
\sin x+\cos x-2\sin x\cos x=0.
\]

Делаем замену:
\[
t=\sin x+\cos x,\qquad t\in[-\sqrt2;\sqrt2].
\]

Тогда:
\[
2\sin x\cos x=t^2-1.
\]

Получаем:
\[
t-(t^2-1)=0.
\]

\[
-t^2+t+1=0.
\]

\[
t^2-t-1=0.
\]

\[
t=\frac{1\pm\sqrt5}{2}.
\]

Проверяем область значений:
\[
\frac{1+\sqrt5}{2}>\sqrt2,
\]
поэтому этот корень не подходит.

Остаётся:
\[
t=\frac{1-\sqrt5}{2}.
\]

Возвращаемся:
\[
\sin x+\cos x=\frac{1-\sqrt5}{2}.
\]

\[
\sqrt2\sin\left(x+\frac{\pi}{4}\right)=\frac{1-\sqrt5}{2}.
\]

\[
\sin\left(x+\frac{\pi}{4}\right)=\frac{1-\sqrt5}{2\sqrt2}.
\]

Это обычное тригонометрическое уравнение.

\sectionline

\section*{\textcolor{darkaccent}{8. Общая оценка выражения \texorpdfstring{$a\sin x+b\cos x$}{a sin x + b cos x}}}

Для любого выражения вида
\[
a\sin x+b\cos x
\]
выполняется оценка:
\[
-\sqrt{a^2+b^2}
\le
a\sin x+b\cos x
\le
\sqrt{a^2+b^2}.
\]

Частный случай:
\[
-\sqrt2\le \sin x+\cos x\le\sqrt2.
\]

\important{Зачем это нужно.}
Так проверяют допустимые значения новой переменной. Если после замены получен корень, не входящий в область значений, его нужно отбросить.

\sectionline

\section*{\textcolor{darkaccent}{9. Типичные ошибки}}

\begin{center}
\renewcommand{\arraystretch}{1.45}
\begin{tabularx}{\linewidth}{@{}p{0.45\linewidth}X@{}}
\toprule
Ошибка & Как правильно \\
\midrule
Пропустить деление общего решения на коэффициент при $x$ &
Если $2x=\frac{\pi}{2}+\pi k$, то $x=\frac{\pi}{4}+\frac{\pi k}{2}$. \\

Поделить на $\cos 2x$ или $\sin x$ без проверки &
Нужно выносить множитель и приравнивать каждый множитель к нулю. \\

Применять формулу двойного угла, хотя она не делает углы одинаковыми &
Если стандартная формула не помогает, ищем формулы суммы, разности или произведения. \\

Сделать замену и забыть область значений &
Для $t=\cos x$: $-1\le t\le1$. Для $t=\sin x+\cos x$: $-\sqrt2\le t\le\sqrt2$. \\

Неверно раскрыть $\sin^4 x$ &
Нужно писать $\sin^4 x=(\sin^2 x)^2=\left(\frac{1-\cos 2x}{2}\right)^2$. \\

Забыть, что $\cos(-x)=\cos x$ &
Косинус чётный: $\cos(-x)=\cos x$. \\
\bottomrule
\end{tabularx}
\end{center}

\newpage

\section*{\textcolor{darkaccent}{Тренировочный блок}}

Решите уравнения. Упражнения расположены по нарастающей сложности.

\begin{enumerate}
\item
\[
\cos x+\cos 5x+\cos 2x=0.
\]

\item
\[
\sin 5x-\sin x=0.
\]

\item
\[
2\sin 3x\cos x-\sin 4x=0.
\]

\item
\[
\sin 7x+\sin 3x-\sin 5x=0.
\]

\item
\[
\cos 4x+\cos 2x-\cos x=0.
\]

\item
\[
\sin^2 x=\frac12.
\]

\item
\[
\cos^2 x-\sin^2 x=\frac12.
\]

\item
\[
\sin^4 x+\sin^4\left(x+\frac{\pi}{4}\right)+\sin^4\left(x-\frac{\pi}{4}\right)=1.
\]

\item
\[
\sin x+\cos x-2\sin x\cos x=0.
\]

\item
\[
2\sin x\cos x+3(\sin x+\cos x)-2=0.
\]
\end{enumerate}

\newpage

\section*{\textcolor{darkaccent}{Ответы}}

Везде $k\in\Z$.

\begin{center}
\renewcommand{\arraystretch}{1.6}
\begin{tabularx}{\linewidth}{@{}cX@{}}
\toprule
№ & Ответ \\
\midrule
1 &
\[
x=\frac{\pi}{4}+\frac{\pi k}{2}
\quad\text{или}\quad
x=\pm\frac{2\pi}{9}+\frac{2\pi k}{3}
\]
\\

2 &
\[
x=\frac{\pi k}{2}
\quad\text{или}\quad
x=\frac{\pi}{3}+\frac{\pi k}{3}
\]
\\

3 &
\[
x=\frac{\pi k}{2}
\]
\\

4 &
\[
x=\frac{\pi k}{2}
\quad\text{или}\quad
x=\pm\frac{\pi}{12}+\frac{\pi k}{2}
\]
\\

5 &
\[
x=\pm\frac{\pi}{9}+\frac{2\pi k}{3}
\quad\text{или}\quad
x=\pm\frac{2\pi}{9}+\frac{2\pi k}{3}
\]
\\

6 &
\[
x=\frac{\pi}{4}+\frac{\pi k}{2}
\quad\text{или}\quad
x=-\frac{\pi}{4}+\frac{\pi k}{2}
\]
\\

7 &
\[
x=\pm\frac{\pi}{6}+\pi k
\]
\\

8 &
\[
x=\frac{k\pi}{2}
\quad\text{или}\quad
x=\pm\frac12\arccos(\sqrt2-1)+\pi k
\]
\\

9 &
\[
\sin\left(x+\frac{\pi}{4}\right)=\frac{1-\sqrt5}{2\sqrt2}
\]
\\

10 &
\[
\sin x+\cos x=\frac{-3+\sqrt{21}}{2}
\]
\\
\bottomrule
\end{tabularx}
\end{center}

\end{document}